\documentclass[english,10pt]{article}

% Encoding and page layout
\usepackage[T1]{fontenc}
\usepackage[utf8]{inputenc}
\usepackage{geometry}
\geometry{verbose,tmargin=3cm,bmargin=3cm,lmargin=2.75cm,rmargin=2.75cm}

% Math
\usepackage{amsmath,amssymb,amsfonts}
\usepackage{mathtools}
\usepackage{mathrsfs}
\usepackage{bm}
\usepackage{bbm}

% Graphics, tables, algorithms
\usepackage{graphicx}
\usepackage{booktabs}
\usepackage{float}
\usepackage{algorithm,algpseudocode}
\newcommand{\algorithmautorefname}{Algorithm}

% Hyperlinks
\usepackage{hyperref}
\hypersetup{
    colorlinks=true,
    linkcolor=blue,
    filecolor=magenta,
    urlcolor=cyan,
    citecolor=red,
}

% Bibliography
\usepackage{natbib}

% Language
\usepackage{babel}
\date{}

% Theorem environments
\usepackage{amsthm}
\numberwithin{equation}{section}

\newtheorem{theorem}{Theorem}[section]
\newtheorem{lemma}{Lemma}[section]
\newtheorem{proposition}{Proposition}[section]
\newtheorem{corollary}{Corollary}[section]
\newtheorem{remark}{Remark}[section]
\newtheorem{assumption}{Assumption}[section]

\title{A Closed-Form Gaussian-Mixture Toy Model \\
       for Validating the Girsanov-Corrected SMC Weight}
\author{Validation note accompanying the main derivation}
\date{}

\begin{document}

\maketitle

\begin{abstract}
This note constructs a one-dimensional Gaussian-mixture toy problem in which the prior,
the noised marginals, the exact score and denoiser, the true posterior, and --
crucially -- the \emph{exact} intermediate likelihood $p(y\mid x_\sigma)$ are all
available in closed form. Because the exact intermediate likelihood is normally
intractable, this toy model lets us simulate both the approximate-guided diffusion
actually used in practice and an exact-guided diffusion that targets the true
posterior exactly, and to check the Girsanov/Doob-transform SMC correction against a
genuine analytical ground truth rather than against itself. We derive the relevant
quantities from scratch, discuss a sign-convention pitfall that this exercise
uncovered in the drift of the reverse-time SDE, and report numerical results:
the corrected sampler converges toward the analytical posterior at the canonical
Monte Carlo rate in the number of particles, up to a small residual bias tied to
the time-discretization step count.
\end{abstract}

\section{The toy model}
\label{sec:model}

We work in one dimension with a $K$-component Gaussian-mixture prior,
\begin{equation}
p(x_0) = \sum_{k=1}^K w_k \, \mathcal N(x_0;\mu_k,\sigma_k^2),
\label{eq:prior}
\end{equation}
noised by the same variance-exploding (VE), EDM-style forward process used
throughout the main derivation \citep{karras2022elucidating}:
\begin{equation}
x_\sigma = x_0 + \sigma\,\epsilon, \qquad \epsilon\sim\mathcal N(0,1),
\label{eq:forward}
\end{equation}
with noise level identified with physical time, $\sigma(t)=t$. We use
$K=2$, $w=(0.5,0.5)$, $\mu=(-3,3)$, $\sigma_k^2=1$ throughout the numerical
experiments below.

\begin{remark}
The conditional law $x_\sigma\mid x_0\sim\mathcal N(x_0,\sigma^2)$ is exactly
Gaussian at \emph{every} finite $\sigma>0$ -- there is nothing asymptotic about
it. The marginal $p_\sigma(x_\sigma)=\int p(x_\sigma\mid x_0)p(x_0)\,dx_0$ is a
Gaussian mixture (below) at every finite $\sigma$ as well; it only becomes
\emph{practically} indistinguishable from a single Gaussian as $\sigma\to\infty$,
once the fixed inter-component separation becomes negligible relative to the
growing noise scale.
\end{remark}

Convolving a Gaussian mixture with independent Gaussian noise leaves the
weights and means untouched and inflates each component's variance:
\begin{equation}
p_\sigma(x) = \sum_k w_k\,\mathcal N(x;\mu_k,\sigma_k^2+\sigma^2).
\label{eq:marginal}
\end{equation}
We adopt a $K$-component Gaussian, linear observation model,
\begin{equation}
y = x_0 + \eta, \qquad \eta\sim\mathcal N(0,r^2).
\label{eq:likelihood}
\end{equation}
A single Gaussian prior would make the Bayes-optimal denoiser affine in $x$,
and the plug-in approximate likelihood introduced below would then be
proportional to the true intermediate likelihood -- same mode, only a scale
difference -- too easy a test of the correction. The mixture prior makes the
denoiser genuinely nonlinear, so the approximation gap has the same
qualitative character it has with a real trained network.

\section{Exact quantities}
\label{sec:exact}

Write $V_k(\sigma):=\sigma_k^2+\sigma^2$ and let
\begin{equation}
\gamma_k(x,\sigma) := \frac{w_k\,\mathcal N(x;\mu_k,V_k(\sigma))}{\sum_j w_j\,\mathcal N(x;\mu_j,V_j(\sigma))}
\label{eq:responsibility}
\end{equation}
be the posterior mixture responsibilities of the noised marginal
\eqref{eq:marginal}. Differentiating $\log p_\sigma$ component-wise gives the
exact score,
\begin{equation}
\nabla_x\log p_\sigma(x) = \sum_k \gamma_k(x,\sigma)\left(-\frac{x-\mu_k}{V_k(\sigma)}\right),
\label{eq:score}
\end{equation}
and the exact Bayes-optimal denoiser $D(x,\sigma):=\mathbb E[x_0\mid x_\sigma=x]$,
\begin{equation}
D(x,\sigma) = \sum_k \gamma_k(x,\sigma)\left[\mu_k + \frac{\sigma_k^2}{V_k(\sigma)}(x-\mu_k)\right],
\label{eq:denoiser}
\end{equation}
consistent with Tweedie's identity $D(x,\sigma)=x+\sigma^2\nabla_x\log p_\sigma(x)$
\citep{efron2011tweedie}. Neither \eqref{eq:score} nor \eqref{eq:denoiser}
requires a trained network -- both replace $D_\theta$ exactly, isolating the
Girsanov/SMC machinery from any network-approximation error.

The posterior $p(x_0\mid y)\propto p(y\mid x_0)p(x_0)$ under
\eqref{eq:likelihood} is again a mixture, with the standard per-component
conjugate update; this is our ground truth.

The intermediate likelihood $p(y\mid x_\sigma)$ -- normally intractable, and
the reason approximate guidance is used in practice at all -- is available
here in closed form. Writing the per-component denoising posterior
$x_0\mid x_\sigma=x,\,k\sim\mathcal N(m_k(x,\sigma),c_k(\sigma))$ with
\begin{equation}
m_k(x,\sigma)=\mu_k+\frac{\sigma_k^2}{V_k(\sigma)}(x-\mu_k), \qquad
c_k(\sigma)=\frac{\sigma_k^2\sigma^2}{V_k(\sigma)},
\label{eq:denoising_posterior}
\end{equation}
and convolving with the observation noise gives
\begin{equation}
p(y\mid x_\sigma=x) = \sum_k \gamma_k(x,\sigma)\,\mathcal N\big(y;\,m_k(x,\sigma),\,r^2+c_k(\sigma)\big).
\label{eq:exact_intermediate}
\end{equation}
The approximate (plug-in) likelihood actually used for guidance in practice
replaces the mixture \eqref{eq:exact_intermediate} with a single Gaussian at
the denoised point estimate,
\begin{equation}
\tilde p_\theta(y\mid x_\sigma=x) = \mathcal N\big(y;\,D(x,\sigma),\,r^2\big).
\label{eq:approx_likelihood}
\end{equation}
Two gaps between \eqref{eq:exact_intermediate} and \eqref{eq:approx_likelihood}
are visible directly from the formulas: the exact version is a genuine mixture
over the per-component means $m_k$, which the plug-in approximation collapses
to a single point; and the exact version's variance is inflated by
$c_k(\sigma)$ -- the residual uncertainty about $x_0$ given only the noisy
$x_\sigma$ -- which the plug-in approximation ignores entirely, silently
treating the denoiser's point estimate as if it were the true $x_0$. That
second gap is exactly what the Girsanov correction exists to fix.

\section{Exact and approximate guided diffusion}
\label{sec:guided}

Let $\tau=T-t$ be the sampling-time reparametrization ($\tau=0$ at pure noise,
$\tau=T$ at clean data) and $\bar a(\tau):=a(T-\tau)$ with
$a(t):=d\sigma(t)^2/dt=2\sigma(t)$ under $\sigma(t)=t$. Anderson's reverse-time
SDE theorem \citep{anderson1982reverse,song2021score} gives the unconditional
reverse process directly in terms of the true score, unambiguously:
\begin{equation}
dZ_\tau = \bar a(\tau)\,\nabla_x\log p_{\sigma(\tau)}(Z_\tau)\,d\tau + \sqrt{\bar a(\tau)}\,dW_\tau,
\qquad \sigma(\tau):=\sigma_{T-\tau}.
\label{eq:unconditional_sde}
\end{equation}
Adding the gradient of a log-likelihood $\ell_\tau$ to the drift steers the
process toward the corresponding posterior:
\begin{equation}
dZ_\tau^\bullet = \bar a(\tau)\Big[\nabla_x\log p_{\sigma(\tau)}(Z_\tau^\bullet) + \nabla_x \ell_\tau(Z_\tau^\bullet)\Big]\,d\tau + \sqrt{\bar a(\tau)}\,dW_\tau.
\label{eq:guided_sde}
\end{equation}
Taking $\ell_\tau(x) := \log p(y\mid x_{\sigma(\tau)}=x)$, the \emph{exact}
intermediate likelihood \eqref{eq:exact_intermediate}, gives the exact-guided
diffusion, whose $\tau=T$ marginal is the true posterior exactly. Taking
$\ell_\tau(x):=\log\tilde p_\theta(y\mid x_{\sigma(\tau)}=x)$, the plug-in
likelihood \eqref{eq:approx_likelihood}, gives the approximate-guided
diffusion actually simulated by a practical guided sampler -- and whose
$\tau=T$ marginal is biased, for the reasons noted above.

\begin{remark}[A sign-convention pitfall]
\label{rem:sign}
It is common in the EDM/denoiser literature \citep{karras2022elucidating} to
write $s_\theta(x,\sigma):=(x-D(x,\sigma))/\sigma^2$, chosen so that the
deterministic probability-flow ODE step comes out in the clean form
$(x-D)/\sigma$. By Tweedie's identity this $s_\theta$ is the \emph{negative}
of the true score, $s_\theta=-\nabla_x\log p_\sigma$. Reverse-SDE formulas of
the form $dX_t=-a(t)s_\theta(X_t,t)\,dt+\ldots$, standard throughout the
score-based generative modeling literature, are only correct as written when
$s_\theta$ denotes the true score directly; substituting the EDM-convention
$s_\theta$ into such a formula introduces a spurious overall sign error in the
unconditional part of the drift. We found exactly this: writing the drift of
\eqref{eq:unconditional_sde} with $+s_\theta$ in place of $\nabla\log p_\sigma$
causes the unguided reverse SDE to diverge (verified numerically -- see
\S\ref{sec:results}), even though the formula is a one-symbol substitution
away from \eqref{eq:unconditional_sde} and each convention is individually
standard in its own home literature. Writing the drift directly in terms of
$\nabla_x\log p_\sigma$ throughout, as in \eqref{eq:unconditional_sde}--\eqref{eq:guided_sde},
avoids the ambiguity entirely and is what we implement.
\end{remark}

Both SDEs are simulated by Euler--Maruyama on a decreasing $\sigma$ grid
(physical time, matching a practical EDM sampler rather than the
increasing-$\tau$ convention used above for the derivation), with step
$\Delta=\sigma_{\mathrm{cur}}-\sigma_{\mathrm{next}}>0$:
\begin{equation}
x_{\mathrm{next}} = x_{\mathrm{cur}} + \bar a(\sigma_{\mathrm{cur}})\Big[\nabla\log p_{\sigma_{\mathrm{cur}}}(x_{\mathrm{cur}}) + \nabla \ell(x_{\mathrm{cur}},\sigma_{\mathrm{cur}})\Big]\Delta + \sqrt{\bar a(\sigma_{\mathrm{cur}})\,\Delta}\;\xi, \qquad \xi\sim\mathcal N(0,1).
\label{eq:euler_maruyama}
\end{equation}

\section{The Girsanov/Doob-transform correction}
\label{sec:girsanov}

Let $P^{\mathrm{pr}}$ and $P^{\mathrm{ap}}$ denote the path measures of
\eqref{eq:unconditional_sde} and of \eqref{eq:guided_sde} with the
\emph{approximate} likelihood, respectively. Itô's formula applied to
$\ell_\tau(Z_\tau^{\mathrm{ap}})$ under $P^{\mathrm{ap}}$ gives
\begin{equation}
d\ell_\tau(Z_\tau^{\mathrm{ap}}) = \Big[\partial_\tau\ell_\tau + b^{\mathrm{ap}}\!\cdot\!\nabla\ell_\tau + \tfrac12\bar a\,\Delta\ell_\tau\Big]d\tau + \sqrt{\bar a}\,\nabla\ell_\tau^\top dW_\tau^{\mathrm{ap}},
\label{eq:ito}
\end{equation}
with $b^{\mathrm{ap}}(x,\tau)=\bar a(\tau)[\nabla\log p_{\sigma(\tau)}(x)+\nabla\ell_\tau(x)]$.
Girsanov's theorem \citep{oksendal2003stochastic} gives, for the drift
difference $\bar a\nabla\ell_\tau$ shared diffusion coefficient $\sqrt{\bar a}$,
\begin{equation}
\frac{dP^{\mathrm{pr}}}{dP^{\mathrm{ap}}} = \exp\left(-\int_0^T\sqrt{\bar a}\,\nabla\ell_\tau^\top dW_\tau^{\mathrm{ap}} - \tfrac12\int_0^T \bar a\,\|\nabla\ell_\tau\|^2\,d\tau\right).
\label{eq:girsanov_toy}
\end{equation}
Substituting the Itô identity \eqref{eq:ito} to eliminate the stochastic
integral, exactly as in the main derivation, gives
\begin{equation}
\frac{dP^{\mathrm{pr}}}{dP^{\mathrm{ap}}} = e^{\ell_0(Z_0)-\ell_T(Z_T)}\exp\left(\int_0^T V_\tau(Z_\tau)\,d\tau\right),
\label{eq:master}
\end{equation}
\begin{equation}
V_\tau(x) := \partial_\tau\ell_\tau(x) + \bar a(\tau)\,\nabla\log p_{\sigma(\tau)}(x)\cdot\nabla\ell_\tau(x) + \tfrac12\bar a(\tau)\|\nabla\ell_\tau(x)\|^2 + \tfrac12\bar a(\tau)\Delta\ell_\tau(x).
\label{eq:V_tau_toy}
\end{equation}
Written directly in terms of the true score, \eqref{eq:V_tau_toy} has no sign
ambiguity; it is the drift-substituted, simplified form of $V_\tau$ used
throughout, and it is what is actually implemented (with $\nabla\log
p_{\sigma(\tau)}=-s_\theta$ if one prefers the EDM-convention $s_\theta$, cf.
Remark~\ref{rem:sign}). Since $\sigma(t)=t$ and $\tau=T-t$, $\partial_\tau\ell_\tau(x) = -\partial_\sigma\ell(x,\sigma)|_{\sigma=\sigma(\tau)}$.

\paragraph{Discretization.} Rather than a naive left-endpoint quadrature of
the whole of \eqref{eq:V_tau_toy}, we use the exact, frozen-state telescoping
identity for the $\partial_\tau\ell_\tau$ term,
$\int_{\tau_{k-1}}^{\tau_k}\partial_\tau\ell_\tau(z_{k-1})\,d\tau = \ell_{\tau_k}(z_{k-1})-\ell_{\tau_{k-1}}(z_{k-1})$,
and a quadrature only for the remaining term
$H(x,\sigma):=\bar a(\sigma)\nabla\log p_\sigma(x)\cdot\nabla\ell(x,\sigma)+\tfrac12\bar a(\sigma)\|\nabla\ell(x,\sigma)\|^2+\tfrac12\bar a(\sigma)\Delta\ell(x,\sigma)$.
The incremental log-weight for a step from $\sigma_{\mathrm{cur}}$ to
$\sigma_{\mathrm{next}}$ at the frozen pre-step state $z$ is
\begin{equation}
\log\hat G = \big[\ell(z;\sigma_{\mathrm{next}}) - \ell(z;\sigma_{\mathrm{cur}})\big] + \int_{\sigma_{\mathrm{next}}}^{\sigma_{\mathrm{cur}}}\! H(z,\sigma)\,d\sigma,
\label{eq:incremental_weight}
\end{equation}
with a one-time correction term $\ell(z_0;\sigma_{\max})$ added before the
first step, reproducing the $e^{\ell_0(Z_0)}$ boundary factor in
\eqref{eq:master} (the discrete analogue of absorbing the initial condition
into the first incremental weight). We combine \eqref{eq:incremental_weight}
with adaptive systematic resampling, triggered whenever the effective sample
size drops below half the particle count.

\section{Numerical results}
\label{sec:results}

We fix $y=0.5$, $r=1$, giving a genuinely bimodal posterior with weights
$(0.182, 0.818)$ on the two components -- an asymmetric but non-degenerate
target. Particles are propagated on the rho-spaced $\sigma$-schedule of
\citep{karras2022elucidating}, $\sigma_{\min}=0.02$, $\sigma_{\max}=20$, with
$1000$ steps.

\paragraph{Histogram comparison.} Figure~\ref{fig:histogram} compares, at
$N=200{,}000$ particles: the exact-guided diffusion (using
\eqref{eq:exact_intermediate}), the uncorrected approximate-guided diffusion
(using \eqref{eq:approx_likelihood} with no reweighting), the SMC-corrected
approximate-guided diffusion (reweighted by \eqref{eq:incremental_weight} and
resampled), and the analytical posterior density. The exact-guided diffusion
matches the analytical posterior to within sampling noise (maximum absolute
density deviation $0.004$, consistent with the Poisson floor at this particle
count), confirming that the base SDE machinery -- score, discretization, drift
sign -- is correct independently of any approximation. The uncorrected
approximate-guided diffusion is visibly biased: narrower, mode-collapsed
toward a single peak (maximum deviation $0.294$, about $67\times$ larger).
The SMC-corrected diffusion recovers the bimodal shape closely, though not
perfectly (maximum deviation in the $0.01$--$0.07$ range depending on run and
bin width) -- the residual gap is discussed below.

\begin{figure}[H]
\centering
\includegraphics[width=0.85\textwidth]{../figures/toy_histogram_comparison.pdf}
\caption{Density comparison at $N=200{,}000$ particles. The exact-guided
diffusion and the SMC-corrected approximate-guided diffusion both recover the
bimodal analytical posterior; the uncorrected approximate-guided diffusion
collapses toward a single mode.}
\label{fig:histogram}
\end{figure}

\paragraph{Convergence in the number of particles.} We estimated the
posterior mean $\mathbb E[x\mid y]$ from the SMC-corrected sampler at
$N\in\{10^3,4\times10^3,1.6\times10^4,6.4\times10^4,2\times10^5\}$, with
$4$--$10$ independent replications per $N$ (fewer at the largest $N$, for
compute reasons). Figure~\ref{fig:mc_error}, left panel, shows both the RMSE
against the true posterior mean ($1.2027$) and the standard deviation across
replications, on a log-log scale, against the canonical $1/\sqrt N$
reference. A least-squares fit of $\log(\text{std})$ against $\log N$ gives
slope $-0.46$, close to the theoretical $-0.5$, confirming the estimator
converges at the usual Monte Carlo rate. The right panel shows the raw bias
(estimate minus truth) with $\pm1$ standard-error bars from the finite
replication count; at every particle count but the largest, the error bar
crosses zero, so the data cannot statistically distinguish a genuine residual
bias from replication noise at that $N$. A parametric fit
$\mathrm{MSE}(N)\approx b^2 + C/N$ under a non-negativity constraint on $b^2$
collapses to $b^2=0$, $C\approx93$ -- i.e., the data as collected is fully
consistent with a purely canonical $C/N$ decay and no detectable bias floor,
though this should not be read as disproving a small bias: a separate
experiment varying the number of discretization \emph{steps} at fixed $N$
found the histogram deviation improving sharply from $250$ to $500$ steps but
then plateauing from $500$ to $2000$ steps, consistent with a small residual
bias tied to the time-discretization order rather than to particle count.
Distinguishing these possibilities conclusively would require substantially
more replications concentrated at large $N$, where the current error bars are
still wide relative to the effect size in question.

\begin{figure}[H]
\centering
\includegraphics[width=0.95\textwidth]{../figures/toy_mc_error_vs_N.pdf}
\caption{Left: RMSE and across-replication standard deviation of the
SMC-corrected posterior-mean estimate versus particle count, log-log scale,
against the canonical $1/\sqrt N$ reference. Right: bias (estimate minus true
posterior mean) with $\pm1$ standard-error bars from the replication count.}
\label{fig:mc_error}
\end{figure}

The same convergence is visible directly in the sampled density itself.
Figure~\ref{fig:histogram_vs_N} overlays the SMC-corrected histogram (a single
representative run, not averaged over replications) against the analytical
posterior at $N=1{,}000$, $16{,}000$, and $200{,}000$ particles. At
$N=1{,}000$ the histogram is visibly jagged, though the bimodal shape is
already present; at $N=16{,}000$ it is smoother but still wobbles around the
analytical curve, with a maximum deviation ($0.036$) that happens to read
slightly \emph{worse} than the $N=1{,}000$ run ($0.023$) on this single-run
metric alone -- consistent with the large run-to-run variance at low $N$
already seen in Figure~\ref{fig:mc_error}, not a reversal of the trend. At
$N=200{,}000$ the histogram sits almost exactly on the analytical curve
(maximum deviation $0.010$). Read together with Figure~\ref{fig:mc_error},
the message is the same one stated quantitatively above: the single-run
histogram error shrinks with $N$ on average, but any individual run at
moderate $N$ can be noisier than a smaller-$N$ run purely by chance, which is
exactly why the particle-count sweep in Figure~\ref{fig:mc_error} was run
with multiple replications rather than read off a single realization.

\begin{figure}[H]
\centering
\includegraphics[width=0.85\textwidth]{../figures/toy_histogram_vs_N.pdf}
\caption{SMC-corrected density histogram (single representative run) at three
particle counts, against the analytical posterior. Each run uses the same
$1000$-step discretization and the same adaptive resampling scheme.}
\label{fig:histogram_vs_N}
\end{figure}

\section{Discussion}

This toy model isolates the Girsanov/Doob-transform SMC correction from
network-approximation error entirely, since the exact score, denoiser,
posterior, and intermediate likelihood are all closed form. It served two
purposes here: it exposed a genuine sign-convention error in the drift of the
reverse-time SDE (Remark~\ref{rem:sign}) before that error could propagate
silently into results on the real PDE problem, and -- once corrected -- it
gave a positive, closed-loop confirmation that the corrected SMC weight
\eqref{eq:incremental_weight} converges to the true posterior at the
canonical Monte Carlo rate in the particle count, up to a small residual tied
to the time-discretization step count rather than to the correction formula
itself. A natural extension, left for future work, is a moderate-dimensional
($d\sim10$--$20$) version of the same construction, still fully closed form,
in which the Laplacian term is a genuine trace over many dimensions and the
Hutchinson estimator used on the real PDE problem can be checked directly
against an exact analytical trace.

\bibliographystyle{plain}
\bibliography{references_2}

\end{document}
